% GENERATED FILE. Edit canonical lessons, then run npm run generate:books.
% canonical-source-hash: fnv1a64:9ffc32f5c5df9119
% canonical-lessons: ZH-C458-qingxu, ZH-C458-xinteng, ZH-C458-xingfen, ZH-C458-qingsong, ZH-C458-fan

\chapter{How I Feel}
\label{ch:zh-how-i-feel}
\hlchaptermodality{458}

\begin{chapteropening}
\textbf{By the end of this chapter:} \emph{I can talk about my mood, feeling for someone, being excited, relaxed or fed up.}

Complete the last lesson of chapter 458: I can talk about my mood, feeling for someone, being excited, relaxed or fed up.
\end{chapteropening}

\section[qíngxù]{qíngxù — emotions; spirits}

\label{lesson:ZH-C458-qingxu}



\begin{quote}
Before the new one: say the Mandarin for to feel deeply grateful, then the Mandarin for a way to solve a problem.
\end{quote}

\subsection*{You'll want to know: qíngxù}
\emph{qíngxù} — \textquotedblleft{}emotions; spirits\textquotedblright{}.

\begin{quote}
\emph{zuìjìn wǒ de qíngxù bú tài hǎo} — My mood hasn't been great lately.
\end{quote}

Its characters come later; for now it is a word for the ear.

\begin{cousinweb}
\emph{qíngxù} is built from \emph{qíng} 'feeling' and \emph{xù} 'the end of a thread', both inherited from Middle Chinese, \emph{dzjeng} and \emph{zjoX} (Guangyun, in Baxter's transcription). The dictionary explains the first, after an older book, as what a person desires, and the second as the loose end of silk; the sense 'mood' is not given there.
\end{cousinweb}

\begin{grammarlens}[title={What you've built}]
One more word for saying how it was.
\end{grammarlens}

\subsection*{Your turn}
\begin{itemize}
\raggedright
  \item \emph{Say it:} \emph{qíngxù}
  \item \emph{Say it:} \emph{qíngxù}, once more
  \item \emph{Say it:} \emph{zuìjìn wǒ de qíngxù bú tài hǎo} — My mood hasn't been great lately.
  \item \emph{Recall:} say the Mandarin for to feel deeply grateful, then the Mandarin for a way to solve a problem, then say \emph{qíngxù} again
\end{itemize}

\begin{tcolorbox}[breakable,colback=teal!4,colframe=teal!35!black,title={Before you move on}]
What does qíngxù mean? (\textquotedblleft{}Emotions; spirits\textquotedblright{}.) Say it once more.
\end{tcolorbox}
\section[\texorpdfstring{\zh{心疼} (xīnténg)}{心疼 (xīnténg)}]{\zh{心疼} (xīnténg) — to feel for someone; to love dearly}

\label{lesson:ZH-C458-xinteng}



\begin{quote}
Before the new one: say the Mandarin for a way to solve a problem, then the Mandarin for one's spirits.
\end{quote}

\subsection*{You'll want to know: \zh{心疼}}
\textbf{\zh{心疼}} — \emph{xīnténg} — \textquotedblleft{}to feel for someone; to love dearly\textquotedblright{}.

\begin{quote}
\textbf{\zh{孩子病了妈妈很心疼}} \emph{háizi bìng le māma hěn xīnténg} — The child is ill, and her mother's heart aches for her.
\end{quote}

\begin{cousinweb}
\textbf{\zh{心疼}} (\emph{xīnténg}) is built from \zh{心} 'heart' and \zh{疼} 'pain', both inherited from Middle Chinese, \emph{sim} and \emph{downg} (Guangyun, in Baxter's transcription), which glosses the second 'pain'. The regular sound change would give \emph{tóng} for the second; why the vowel differs is not settled.
\end{cousinweb}

\begin{grammarlens}[title={What you've built}]
One more word for saying how it was.
\end{grammarlens}

\subsection*{Your turn}
\begin{itemize}
\raggedright
  \item \emph{Say it:} \emph{xīnténg}
  \item \emph{Say it:} \emph{xīnténg}, once more
  \item \emph{Say it:} \emph{háizi bìng le māma hěn xīnténg} — The child is ill, and her mother's heart aches for her.
  \item \emph{Recall:} say the Mandarin for a way to solve a problem, then the Mandarin for one's spirits, then say \emph{xīnténg} again
\end{itemize}

\begin{tcolorbox}[breakable,colback=teal!4,colframe=teal!35!black,title={Before you move on}]
What does \zh{心疼} mean? (\textquotedblleft{}To feel for someone; to love dearly\textquotedblright{}.) Say it once more.
\end{tcolorbox}
\section[xīngfèn]{xīngfèn — thrilled; keyed up}

\label{lesson:ZH-C458-xingfen}



\begin{quote}
Before the new one: say the Mandarin for one's spirits, then the Mandarin for to feel for someone.
\end{quote}

\subsection*{You'll want to know: xīngfèn}
\emph{xīngfèn} — \textquotedblleft{}thrilled; keyed up\textquotedblright{}.

\begin{quote}
\emph{míngtiān qù Zhōngguó wǒ gǎnjué hěn xīngfèn} — I'm going to China tomorrow, and I feel really excited.
\end{quote}

Its characters come later; for now it is a word for the ear.

\begin{cousinweb}
\emph{xīngfèn} is built from \emph{xīng} 'to rise' and \emph{fèn} 'to spread the wings', both inherited from Middle Chinese, \emph{xing} and \emph{pjunH} (Guangyun, in Baxter's transcription), which glosses the first 'to rise' and 'to flourish' and the second as a bird spreading its feathers. The sense 'excited' is not given there.
\end{cousinweb}

\begin{grammarlens}[title={What you've built}]
One more word for saying how it was.
\end{grammarlens}

\subsection*{Your turn}
\begin{itemize}
\raggedright
  \item \emph{Say it:} \emph{xīngfèn}
  \item \emph{Say it:} \emph{xīngfèn}, once more
  \item \emph{Say it:} \emph{míngtiān qù Zhōngguó wǒ gǎnjué hěn xīngfèn} — I'm going to China tomorrow, and I feel really excited.
  \item \emph{Recall:} say the Mandarin for one's spirits, then the Mandarin for to feel for someone, then say \emph{xīngfèn} again
\end{itemize}

\begin{tcolorbox}[breakable,colback=teal!4,colframe=teal!35!black,title={Before you move on}]
What does xīngfèn mean? (\textquotedblleft{}Thrilled; keyed up\textquotedblright{}.) Say it once more.
\end{tcolorbox}
\section[qīngsōng]{qīngsōng — relaxed; light and easy}

\label{lesson:ZH-C458-qingsong}



\begin{quote}
Before the new one: say the Mandarin for to feel for someone, then the Mandarin for thrilled.
\end{quote}

\subsection*{You'll want to know: qīngsōng}
\emph{qīngsōng} — \textquotedblleft{}relaxed; light and easy\textquotedblright{}.

\begin{quote}
\emph{kǎoshì yǐhòu wǒ gǎnjué hěn qīngsōng} — After the exam I feel really relaxed.
\end{quote}

Its characters come later; for now it is a word for the ear.

\begin{cousinweb}
\emph{qīngsōng} is built from \emph{qīng} 'light' and \emph{sōng} 'loose', both inherited from Middle Chinese, \emph{khjieng} and \emph{sowng} (Guangyun, in Baxter's transcription). The dictionary glosses the first as light against heavy, and lists the second under the spelling for loose, tousled hair; the sense 'relaxed' is not given there.
\end{cousinweb}

\begin{grammarlens}[title={What you've built}]
One more word for saying how it was.
\end{grammarlens}

\subsection*{Your turn}
\begin{itemize}
\raggedright
  \item \emph{Say it:} \emph{qīngsōng}
  \item \emph{Say it:} \emph{qīngsōng}, once more
  \item \emph{Say it:} \emph{kǎoshì yǐhòu wǒ gǎnjué hěn qīngsōng} — After the exam I feel really relaxed.
  \item \emph{Recall:} say the Mandarin for to feel for someone, then the Mandarin for thrilled, then say \emph{qīngsōng} again
\end{itemize}

\begin{tcolorbox}[breakable,colback=teal!4,colframe=teal!35!black,title={Before you move on}]
What does qīngsōng mean? (\textquotedblleft{}Relaxed; light and easy\textquotedblright{}.) Say it once more.
\end{tcolorbox}
\section[fán]{fán — annoyed; fed up}

\label{lesson:ZH-C458-fan}



\begin{quote}
Before the new one: say the Mandarin for thrilled, then the Mandarin for relaxed and at ease.
\end{quote}

\subsection*{You'll want to know: fán}
\emph{fán} — \textquotedblleft{}annoyed; fed up\textquotedblright{}.

\begin{quote}
\emph{zhǎo bu dào wǒ de xínglixiāng wǒ hěn fán} — I can't find my suitcase, and I'm really fed up.
\end{quote}

Its characters come later; for now it is a word for the ear.

\begin{cousinweb}
\emph{fán} is inherited from Middle Chinese \emph{bjon} (Guangyun, in Baxter's transcription), glossed 'weary', and explained after an older book as a hot, aching head. The same syllable is the first half of \emph{máfan} 'trouble'.
\end{cousinweb}

\begin{grammarlens}[title={What you've built}]
One more word for saying how it was.
\end{grammarlens}

\subsection*{Your turn}
\begin{itemize}
\raggedright
  \item \emph{Say it:} \emph{fán}
  \item \emph{Say it:} \emph{fán}, once more
  \item \emph{Say it:} \emph{zhǎo bu dào wǒ de xínglixiāng wǒ hěn fán} — I can't find my suitcase, and I'm really fed up.
  \item \emph{Recall:} say the Mandarin for thrilled, then the Mandarin for relaxed and at ease, then say \emph{fán} again
\end{itemize}

\begin{tcolorbox}[breakable,colback=teal!4,colframe=teal!35!black,title={Before you move on}]
What does fán mean? (\textquotedblleft{}Annoyed; fed up\textquotedblright{}.) Say it once more.
\end{tcolorbox}
